\section{Theoretical Analysis}
\label{sec:theory}

We provide formal justification for why relation-conditioned uncertainty is necessary and sufficient for novel context detection.

\subsection{Impossibility of Relation-Agnostic Detection}

\begin{definition}[Novel Context]
A test triple $(h, r, t)$ is in \textbf{novel context} if $\cov(h, r) = 0$ or $\cov(t, r) = 0$.
\end{definition}

\begin{theorem}[Impossibility]
\label{thm:impossibility}
Let $U: \mathcal{E} \to \mathbb{R}$ be any uncertainty measure that depends only on the entity, not the queried relation. Then for the task of distinguishing novel context triples from in-distribution triples:
\[
\text{AUROC}(U) \leq 0.5 + \epsilon
\]
for arbitrarily small $\epsilon > 0$.
\end{theorem}

\begin{proof}
Consider an entity $e$ that has been observed with relations $r_1, r_2, \ldots, r_m$ but not with relations $r_{m+1}, \ldots, r_n$.

For any relation-agnostic uncertainty $U(e)$:
\begin{itemize}
    \item Query $(e, r_1, ?)$ has uncertainty $U(e)$ and is in-distribution
    \item Query $(e, r_{m+1}, ?)$ has uncertainty $U(e)$ and is novel context
\end{itemize}

Since both queries have identical uncertainty $U(e)$, no threshold can separate them. Across the test set, the AUROC for distinguishing these cases approaches 0.5 (random guessing).
\end{proof}

\paragraph{Empirical Validation.}
We verify Theorem~\ref{thm:impossibility} by training a UKGE-style model with entity-level variance only. On FB15k-237, this achieves AUROC = 0.50 for novel context detection, confirming the theoretical prediction.

\subsection{Sufficiency of Relation-Conditioned Variance}

\begin{theorem}[Sufficiency]
\label{thm:sufficiency}
Let $U(e, r) = \sigma^2(e | r) \cdot (1 + k \cdot (1 - \cov(e, r)))$ for some $k > 0$. Then:
\[
\lim_{k \to \infty} \text{AUROC}(U) = 1
\]
\end{theorem}

\begin{proof}
For in-distribution queries where $\cov(e, r) = 1$:
\[
U_\text{ID} = \sigma^2(e | r) \cdot 1 = \sigma^2(e | r)
\]

For novel context queries where $\cov(e, r) = 0$:
\[
U_\text{OOD} = \sigma^2(e | r) \cdot (1 + k)
\]

The ratio is:
\[
\frac{U_\text{OOD}}{U_\text{ID}} = 1 + k
\]

As $k \to \infty$, $U_\text{OOD} \gg U_\text{ID}$, yielding perfect separation.
\end{proof}

\paragraph{Practical Implications.}
In practice, we do not need $k \to \infty$. Our experiments show:
\begin{itemize}
    \item $k = 2$ (boost $= 3\times$): AUROC $\approx 0.97$
    \item $k = 4$ (boost $= 5\times$): AUROC $\approx 0.999$
    \item Learnable $k$ converges to $\approx 2.5$ (boost $\approx 3.5\times$)
\end{itemize}

The learned value balances novel context detection with maintaining meaningful uncertainty gradations for in-distribution queries.

\subsection{Why Coverage Cannot Be Learned}

One might ask: can a model learn to track coverage implicitly, without explicit supervision?

\paragraph{Empirical Evidence.}
We trained models with coverage-aware regularization, encouraging high uncertainty for held-out entity-relation pairs. Result: AUROC \emph{decreased} from 0.63 to 0.58. The model memorizes specific pairs rather than learning the general pattern.

\paragraph{Theoretical Intuition.}
Coverage is a \emph{structural} property of the training data, orthogonal to the \emph{semantic} patterns captured by embeddings. Neural networks learn to interpolate semantic patterns but cannot reliably extrapolate structural absence (what was \emph{not} observed).

This motivates our explicit coverage tracking, which provides a guaranteed signal that learned components cannot discover autonomously.
